\rtitle{Evolution of Traffic Flow Visualization in Civil Engineering Education}
\rnote{The concise run made more tool calls than the detail run (122 against 74) but keeps only four sources and three figures: one figure per idea, and a single paragraph each for simulation and dashboards.}
Traffic flow theory and highway capacity analysis have shifted from static diagrams to dynamic, data-driven visualizations. Early education relied on \textbf{\clue{time-space diagrams}} to show vehicle trajectories and \textbf{\clue{flow-density relationships}} to explain congestion thresholds.

\rtrace{turn 8 \texttt{open(url)} \(\rightarrow\) turn 9 \texttt{open(id=page\_img\_49)}}{Source page: \emph{Traffic flow - Wikipedia}; caption on the page: ``Figure 1. Time Space diagram''\par\textbf{Why it was used.} Same Wikipedia figure as in the detail run; here one figure carries both the speed and the gap definitions.}
\rimagebox{figures/showcase/traffic/concise/page_img_49__Time_Space_Diag_Figure_1__e01e89d919.jpg}{Time-space diagram showing vehicle trajectories with labeled speed and headway}{\rbox{0.4}{0.43}{0.75}{0.53}{}\rbox{0.35}{0.27}{0.43}{0.5}{}\rbox{0.72}{0.71}{0.8}{0.82}{}}
The time-space diagram plots position versus time, with each curve representing a vehicle's path. \fact{The slope indicates speed, while vertical spacing shows headway} \rcite{1}. These diagrams remain foundational for teaching shockwave theory and bottleneck analysis.

\rtrace{turn 8 \texttt{open(url)} \(\rightarrow\) turn 9 \texttt{open(id=page\_img\_51)}}{Source page: \emph{Traffic flow - Wikipedia}; caption on the page: ``Figure 3. Flow Density relationship''\par\textbf{Why it was used.} The single figure the concise report uses for congestion thresholds: critical and jam density.}
\rimagebox{figures/showcase/traffic/concise/page_img_51__960px-Flow_Density_Relationship__57ffd32caf.jpg}{Flow-density relationship showing free flow, critical density, and jam density}{\rbox{0.14}{0.18}{0.44}{0.24}{}\rbox{0.68}{0.18}{0.93}{0.24}{}}
The flow-density curve illustrates the transition from free-flow to congested conditions. The peak at critical density represents maximum throughput, while the downward slope toward jam density explains why adding vehicles can reduce flow \rcite{2}.

Modern education now integrates \textbf{\clue{microsimulation}} and \textbf{\clue{real-time dashboards}}. Software like TransModeler converts \clue{Highway Capacity Manual (HCM)} analyses into animated 3D visualizations, letting students observe intersection operations and level-of-service impacts dynamically \rcite{3}.

\rtrace{turn 2 \texttt{open(url)} \(\rightarrow\) turn 5 \texttt{open(id=page\_img\_3)}}{Source page: \emph{TransModeler HCM Compatibility}; caption on the page: ``TransModeler traffic simulation showing level of service and traffic impact for signalized intersections and alternative intersections''}
\rimagebox{figures/showcase/traffic/concise/page_img_3__transmodeler-highway-capacity-software__ff22233dda.jpg}{TransModeler simulation showing signalized intersections with level-of-service analysis}{}
Contemporary research also demonstrates how \clue{connected vehicle trajectory data} can be visualized for corridor-level performance assessment. A Purdue study developed a framework that translates millions of GPS points into actionable signal performance measures, including \fact{arrivals on green and downstream blockage} \rcite{4}. This progression---from hand-drawn diagrams to interactive simulations---reflects the field's broader move toward data-rich, visualization-centered pedagogy.

\rsec{References}
\begin{rlist}
\item Traffic flow - Wikipedia\rurl{https://en.wikipedia.org/wiki/Traffic_flow}
\item Level of service (transportation) - Wikipedia\rurl{https://en.wikipedia.org/wiki/Level_of_service_(transportation)}
\item TransModeler HCM Compatibility\rurl{https://www.caliper.com/transmodeler/hcm-compatible.htm}
\item Trajectory-Based Arterial Traffic Signal Performance Measures Reports\rurl{https://docs.lib.purdue.edu/jtrpaffdocs/45/}
\end{rlist}
\rgate{\item Submission 1 (turn 24): {\color{rframe}\textbf{WARN}} --- 1,403 chars, 4 images, 4 sources; \texttt{reference\_unreachable}
\item Submission 2 (turn 25): {\color{tok}\textbf{OK}} --- 1,540 chars, 3 images, 4 sources}
